%% kapa: exam=1 round=34 year=2023 period=2 form=A subject=accounting subject_ko=회계학 seq=36 number=76
%% source: Q-Net 감정평가사 기출문제 (cst00309) artlSeq=5211789 file=raw/1008/5211789/2023년도 제34회 감정평가사 1차시험 2교시.pdf
%% conversion: kordoc 4.24.51 -> tools/gen1.py (auto); review=proofread-sonnet; flags=none
\begin{question}{76}
당기에 설립된 (주)감평은 결합공정을 통하여 제품 X와 Y를 생산·판매한다. 제품 X는 분리점에서 즉시 판매하고 있으나, 제품 Y는 추가가공을 거쳐 판매한다. 결합원가는 균등이익률법에 의해 각 제품에 배분되며, 직접재료는 결합공정 초에 전량 투입되고 전환원가는 결합공정 전반에 걸쳐 균등하게 발생한다. 당기에 (주)감평은 직접재료 3,000단위를 투입하여 2,400단위를 제품으로 완성하고, 600단위는 기말재공품(전환원가 완성도 50\%)으로 남아 있다. 당기에 발생한 직접재료원가와 전환원가는 각각 ￦180,000과 ￦108,000이다. (주)감평의 당기 생산 및 판매 관련 자료는 다음과 같다.
\par{\small\centering\begin{tabular}{|c|c|c|c|c|}\hline
구 분 & 생산량 & 판매량 & \makecell{단위당\\추가가공원가} & \makecell{단위당\\판매가격} \\ \hline
제품 X & 800단위 & 800단위 & － & ￦150 \\ \hline
제품 Y & 1,600 & 900 & ￦15 & 200 \\ \hline
\end{tabular}\par}
제품 Y의 단위당 제조원가는? (단, 공손 및 감손은 발생하지 않는다.)
\begin{choices}
  \item ￦100
  \item ￦105
  \item ￦110
  \item ￦115
  \item ￦120
\end{choices}
\end{question}
\answer{5}
